\section{Boundary cases and detailed proof edges}
\label{app:proofs}

This appendix expands the combinatorial and stability edges most vulnerable to
off-by-one or hidden-decision assumptions.  All indices below are local indices
in the current eligible knot list; mapping them back to original sample indices
preserves order and nesting.

\begin{lemma}[Progress of a nonterminal centred interval]
\label{lem:progress}
Suppose the current interval starts at eligible index $r$ and the centred walk
closes it at $t$ before the terminal endpoint.  Then $t-r\ge2$.
\end{lemma}
\begin{proof}
Let $p\ge r+1$ be the first active curvature after any insignificant entries.
A nonterminal closure requires a later active curvature $q>p$ of the opposite
sign.  There are three possible selected locations.  The first opposite point
is $q\ge p+1\ge r+2$.  A single eligible zero between the two signs is selected
at one index after the last active point, again at least $p+1$.  On a longer
zero plateau, the last point on the departing side is eligible only when it is
not the first active point; hence it is at least $p+1$.  Otherwise the arriving
side is selected.  The implementation's final spacing guard enforces the same
lower bound in degenerate branches.  Thus every nonterminal case has
$t-r\ge2$.
\end{proof}

\begin{lemma}[Halving recurrence]
\label{lem:halving}
If one level starts with $N\ge1$ eligible intervals and produces $N'$ chord
intervals, then $N'\le\lceil N/2\rceil$.
\end{lemma}
\begin{proof}
The produced intervals are consecutive and partition the $N$ old intervals.
By Lemma~\ref{lem:progress}, every produced interval except possibly the last
contains at least two old intervals.  The last contains at least one.  Hence
$N\ge2(N'-1)+1=2N'-1$, so $N'\le(N+1)/2$ and integrality gives the claim.  If
the last interval also contains at least two, the stronger
$N'\le\lfloor N/2\rfloor$ holds.
\end{proof}

\begin{proof}[Detailed completion of Theorem~\ref{thm:finite} and Corollary~\ref{cor:sparse}]
Every new knot is selected from the current eligible list, so the mapped knot
sets are nested.  Let $N_0=n-1$ and let $N_j$ be the number of intervals after
level $j$.  Lemma~\ref{lem:halving} and the identity
$\lceil\lceil a/2\rceil/2\rceil=\lceil a/4\rceil$ give inductively
$N_j\le\lceil N_0/2^j\rceil$.  For $n\ge3$, choosing
$j=\lceil\log_2(n-1)\rceil$ makes $N_j=1$, so the endpoint chord has been
recorded and the implementation stops.  For $n=2$, it records that chord in
one level.  This proves the stated
$\max\{1,\lceil\log_2(n-1)\rceil\}$ bound.  If a sparse implementation spends
$O(N_j)$ work at level $j$, the geometric recurrence gives $O(n)$ total work.
More explicitly, for integer $N_0\ge1$,
$\lceil N_0/2^j\rceil\le(N_0-1)/2^j+1$, so
$\sum_{j=1}^J N_j<N_0-1+J\le2N_0$.  Adding one endpoint per stored knot set
gives $\sum_{j=1}^J|K_j|<2N_0+J$, the stated storage bound.  This is the
complexity of the implemented knot-only representation; materializing every
length-$n$ level remains $O(n\log n)$ in the worst case.
\end{proof}

\begin{lemma}[Centred decision margin]
\label{lem:decision}
On a uniform grid of spacing $h$, define $\gamma$ and $\eta$ as in
Theorem~\ref{thm:local-stability}.
If
$\|e\|_\infty<\min\{\gamma h/4,\eta h/8\}$, every curvature sign and every
adjacent opposite-sign magnitude comparison made by the centred walk is
unchanged.
\end{lemma}
\begin{proof}
The divided-curvature perturbation stencil is
$(1,-2,1)/h$, so every coordinate changes by at most
$4\|e\|_\infty/h$.  The first strict inequality prevents zero crossing.  For
two curvature coordinates $a,b$, the reverse triangle inequality gives
\[
\left|(|a+\delta a|-|b+\delta b|)-(|a|-|b|)\right|
\le |\delta a|+|\delta b|
\le8\|e\|_\infty/h.
\]
The second strict inequality therefore preserves every comparison outcome.
Once signs are fixed there are no zero-curvature branches, and the potential
comparison pairs are the adjacent entries across sign-run boundaries.  These
data determine the same knot walk.
\end{proof}

\begin{proof}[Detailed completion of Theorem~\ref{thm:local-stability}]
Lemma~\ref{lem:decision} fixes the knot set.  At any sample between two fixed
knots, the interpolation error is a convex combination of the two endpoint
errors and is therefore at most $\|e\|_\infty$.  This proves the baseline
bound.  The detail perturbation is $e-[T(y+e)-Ty]$, so the triangle inequality
gives the factor two.  If $\eta=0$, no positive centred-rule radius follows:
the explicit $(1,1,-1)$ counterexample in the main text changes the selected
side while preserving every sign.
\end{proof}

\begin{lemma}[Essential class in the signed persistence metric]
\label{lem:essential}
For a nonempty connected finite path, the ordinary zero-dimensional superlevel
diagram has exactly one essential class.  Under a finite-cost bottleneck
matching between two functions on that path, the essential classes match each
other; the remaining matching restricts to their finite diagrams.
\end{lemma}
\begin{proof}
At the bottom of the filtration the entire path is present and connected, so
exactly one component survives indefinitely.  A diagram point with an infinite
death coordinate has infinite distance from the diagonal and from every finite
point.  Therefore any finite-cost matching pairs the two essential points.
Their cost is the absolute difference of their birth heights, which are the
two function maxima.  Removing that pair leaves precisely a matching of the
finite parts.
\end{proof}

\begin{proof}[Detailed completion of Theorem~\ref{thm:persistence}]
The $i$th divided-curvature row is
\[
\left(h_{i-1}^{-1},
 -(h_{i-1}^{-1}+h_i^{-1}),h_i^{-1}\right),
\]
whose absolute row sum is $2(h_{i-1}^{-1}+h_i^{-1})$.  Taking the maximum gives
the exact induced $\ell_\infty$ norm $C_x$.  Coordinatewise positive part is
1-Lipschitz, so both signed vertex functions differ by at most
$C_x\|y-z\|_\infty$.  Their piecewise-linear extensions over a common edge
differ by a convex combination of endpoint differences and obey the same bound.
Apply persistence-diagram stability to their negatives, which converts
superlevel filtrations to sublevel filtrations without changing the sup norm.
Lemma~\ref{lem:essential} identifies the essential and finite costs used in
$d_{\rm SC}$.  Taking the maximum over signs proves the metric bound.  Finally,
a finite bar $(b,d)$ has diagonal cost $(b-d)/2$; if its lifetime is strictly
greater than twice the metric bound, a bottleneck matching attaining that bound
cannot send it to the diagonal.
\end{proof}

\begin{proof}[Domain detail for Theorem~\ref{thm:prox}]
For $n\ge3$, the divided-curvature matrix has rank $n-2$: given any target
curvature vector, choose an initial value and slope and integrate the target
curvatures successively to construct a preimage.  Thus $D$ is surjective.
Since a proper $\phi$ has nonempty domain, $\phi\circ D$ is proper; it is closed
and convex by continuity and linearity of $D$.  Adding
$\|y-z\|_2^2/2$ makes the objective coercive and strongly convex, so its
minimizer exists and is unique.  The standard proximal theorem now applies to
$g=\phi\circ D$.  If $P=\operatorname{prox}_g$ is firmly nonexpansive, then so
is $I-P$: writing $d=y-z$ and $p=Py-Pz$, firm nonexpansiveness gives
$\|p\|^2\le\langle p,d\rangle$, which directly implies
\[
\|(I-P)y-(I-P)z\|^2\le
\langle (I-P)y-(I-P)z,y-z\rangle.
\]
Finally, $D\ell=0$ gives
$g(w+\ell)=g(w)$; translating the optimization variable establishes affine
equivariance.  The $n=2$ case has the zero curvature map and is immediate.
\end{proof}
